\documentclass[12pt]{article}
\usepackage{graphicx} % Required for inserting images
\usepackage[margin=0.5in]{geometry}
\usepackage{amsmath, amssymb}
\usepackage{mathrsfs}


\usepackage{soul}


% Dr. David Brown Commands
\newcommand{\ds}{\displaystyle}
\newcommand{\R}{\mathbb{R}}
\newcommand{\N}{\mathbb{N}}
\newcommand{\Q}{\mathbb{Q}}
\newcommand{\C}{\mathbb{C}}


% Commands to get script r
\usepackage{calligra}
\DeclareMathAlphabet{\mathcalligra}{T1}{calligra}{m}{n}
\DeclareFontShape{T1}{calligra}{m}{n}{<->s*[2.2]callig15}{}
\newcommand{\scripty}[1]{\ensuremath{\mathcalligra{#1}}}

% Unit commands
\newcommand{\degree}{\text{°}}
\newcommand{\unitSpeed}{\frac{\text{m}}{\text{s}}}
\newcommand{\unitAcceleration}{\frac{\text{m}}{\text{s}^2}}

% Constant commands
\newcommand{\avogadro}{6.022\times10^{23}}

\newcommand{\boltzmannConstK}{1.381\times10^{-23}\frac{\text{J}}{\text{K}}}
\newcommand{\boltzmannConstInEV}{8.617\times10^{-5}\frac{\text{eV}}{\text{K}}}
\newcommand{\planckConst}{6.626\times10^{-34}\text{J}\cdot\text{s}}

\newcommand{\atmP}{1.01325\times10^5\,\text{Pa}}

% Commands to easily write partial derivatives
\newcommand{\partDeriv}[2]{\frac{\partial #1}{\partial #2}}
\newcommand{\doublePartDeriv}[2]{\frac{\partial^2 #1}{\partial^2 #2}}


\title{Problem 7.18}
\author{Gabriel Decker}


\begin{document}


\maketitle
\begin{flushleft}

In this problem, we will consider a type of particle where zero, one, or two particles can share an energy state.
We will derive the occupancy, or the average number of particles in the state.
This is simply an average value problem involving probability.
\begin{equation}
    \overline{n}=\sum_nn\mathcal{P}(n)
\end{equation}
where $\mathcal{P}(n)$, the probability of a number of particles being in a state is the following.
\begin{equation}
    \mathcal{P}(n)=\frac{1}{\mathcal{Z}}e^{-n(\epsilon-\mu)/kT}
\end{equation}
where $\mathcal{Z}$, or the grand partition function for this system, is the sum of the Gibbs factors.
In this case, we're considering that the energy of the state is $E=n\epsilon$.
This gives us Gibbs factors of $e^{-n(\epsilon-\mu)/kT}$.\\~\\

First, let's determine the grand partition function.
It is the sum of possible states for our energy level, which depends on $n$.
Therefore, since for our hypothetical particle, $n=0,1,2$, this implies the following.
$$\mathcal{Z}=\sum_{n=0}^2e^{-n(\epsilon-\mu)/kT}$$
$$\implies\mathcal{Z}=e^{-0(\epsilon-\mu)/kT}+e^{-1(\epsilon-\mu)/kT}+e^{-2(\epsilon-\mu)/kT}$$
$$\implies\mathcal{Z}=1+e^{-(\epsilon-\mu)/kT}+\left(e^{-(\epsilon-\mu)/kT}\right)^2$$\\~\\

We have the grand partition function, so we can now calculate the probability of $n$ particles sharing an energy state, again, the possible values for $n$ are $n=0,1,2$.
$$\mathcal{P}(n)=\frac{1}{\mathcal{Z}}e^{-n(\epsilon-\mu)/kT}$$
$$\implies\mathcal{P}(n)=\frac{\ds e^{-n(\epsilon-\mu)/kT}}{\ds1+e^{-(\epsilon-\mu)/kT}+\left(e^{-(\epsilon-\mu)/kT}\right)^2}$$\\~\\

Now that we have the probability function, we can finally get the occupancy distribution.
$$\overline{n}=\sum_nn\mathcal{P}(n)$$
$$\implies\overline{n}=\sum_{n=0}^2n\frac{\ds e^{-n(\epsilon-\mu)/kT}}{\ds1+e^{-(\epsilon-\mu)/kT}+\left(e^{-(\epsilon-\mu)/kT}\right)^2}$$
$$\implies\overline{n}=0+1\frac{\ds e^{-1(\epsilon-\mu)/kT}}{\ds1+e^{-(\epsilon-\mu)/kT}+\left(e^{-(\epsilon-\mu)/kT}\right)^2}+2\frac{\ds e^{-2(\epsilon-\mu)/kT}}{\ds1+e^{-(\epsilon-\mu)/kT}+\left(e^{-(\epsilon-\mu)/kT}\right)^2}$$
$$\implies\overline{n}=\frac{1}{\ds e^{(\epsilon-\mu)/kT}+1+e^{-(\epsilon-\mu)/kT}}+2\frac{1}{\ds \left(e^{(\epsilon-\mu)/kT}\right)^2+e^{(\epsilon-\mu)/kT}+1}$$
We can define $x=(\epsilon-\mu)/kT$ to simplify the equation.
$$\implies\overline{n}=\frac{1}{\ds e^{x}+1+e^{-x}}+\frac{2}{\ds \left(e^{x}\right)^2+e^{x}+1}$$\\~\\

With this, we can plot the average number of particles as a function of $\epsilon$ for a number of $T$.\\
\includegraphics[width=0.8\linewidth]{energy_vs_occupancy.png}





\end{flushleft}

\end{document}
